% Chapter 15: Direct Interface
% Scene function: The protocol demands Tert show himself. 37 minutes
% of obsessive positioning. Gordon talks to him. Brief, awkward
% conversation. The craftsman's weak point, fully exposed.

\chapter{Direct Interface}

The mirror was not enough.

I know this because the Phase 4 checklist specifies a minimum of two direct-perception events per week, and the mirror counts as one. The checklist does not distinguish between a quarter-second peripheral reflection and a full-duration visual confrontation. Both check the same box. But Prim reads the forms, and Prim will notice if week four has only one event where the protocol calls for two.

I have never been good at this part.

Atmospheric work is what I do. It is what I have always done. The slow build, the sub-threshold erosion, the wrongness you cannot locate. These require patience and technical precision and a tolerance for lying on concrete floors at three in the morning. They do not require being looked at.

Direct interface requires being looked at.

Apparitions is an entire department built around being looked at. They train for it. They workshop their faces. They have opinions about jawlines. I attended one of their sessions once. I left during a debate about eye contact duration. The session assumed that being perceived by the subject was the goal, the skill, the art form. I found this baffling then and I find it baffling now. The entire craft of atmospheric haunting is built on the principle that what you cannot see is worse than what you can. Showing yourself is an admission that the atmosphere was not enough.

But the protocol calls for it, and the form has a field for it.

Friday. 11:20 p.m. Gordon is in bed, light off, not yet asleep. I am in the hallway outside his bedroom. The door is open most of the way. The hallway is dark.

I need to stand in the doorway. That is all. Stand in the doorway, hold position, allow the subject to achieve visual contact. Three seconds minimum per protocol. I have done this before, on other cases, and it has always been the worst three seconds of the assignment.

I have already swept the house for cameras. This was the first thing I did on the first visit. It is always the first thing. Modern subjects install them everywhere: doorbell cameras, motion sensors, devices that record continuously to a server I cannot reach. Gordon does not have any of these. Gordon has a landline and a laptop he closes when he is not using it. I am grateful for this in a way that would have been unnecessary thirty years ago, when the primary surveillance risk was a nosy neighbor with binoculars and a drinking problem.

I step into the doorway.

Gordon does not open his eyes.

His breathing is steady but not deep. He is in the space between awake and asleep where the body has decided to rest and the mind has not agreed. The protocol requires visual contact. I cannot get visual contact with a man whose eyes are closed. I hold position.

Five minutes. I adjust two inches to the left. The streetlight from the living room casts a faint line along the hallway floor, and at this angle it will backlight my silhouette without illuminating detail. This is the correct choice. I hold the position. I consider it. I adjust one inch to the right.

Twelve minutes. The shadow is wrong. At this distance from the door frame, the shadow falls behind me rather than beside me, which means Gordon will see a shape without a shadow, which is either more effective or less effective depending on which school of thought you subscribe to. The Apparitions workshop spent forty minutes on shadow theory. I left before the conclusion. I adjust back to the left.

Twenty minutes. I consider my posture. I am standing the way I normally stand, which is the way a person stands when they are thinking about joist intervals and not about being perceived. This is presumably not the correct posture for a direct-perception event. I straighten slightly. I relax slightly. I cannot tell which is more unsettling.

Twenty-nine minutes. Gordon shifts. He does not wake. I have been doing this for twenty years and I do not know what I look like.

Thirty-seven minutes. Gordon opens his eyes. He looks at me.

``Who is there?'' he says.

I do not move. The protocol says to withdraw after visual contact. I should be gone already.

``I can see you,'' he says. He sits up. ``What are you doing in my house?''

His voice is not frightened. It is the voice of a man who spent thirty years managing classrooms and does not panic when someone is where they should not be. There is authority in it. I am being addressed the way a student is addressed when they have been caught somewhere they do not belong.

I should step back. I should be in the hallway. I should be gone.

``Wrong house,'' I say.

I do not know why I say this. It is not in the protocol. There is no protocol for speaking to the subject. The Operations Manual does not cover this. Nobody has ever briefed me on what to say when a man you have been tormenting for a month asks you what you are doing in his house and you do not have an answer that fits on a form.

``The back door was open,'' I say. ``I apologize.''

Gordon looks at me for a long moment. Then his voice changes. The authority goes out of it. What replaces it is worse.

``Are you all right?'' he says.

I am gone before he can turn on the light.

Gordon gets out of bed. He checks every room in the house. He checks the front door. He checks the back door. He does this methodically, the way a person checks a house when they believe a confused stranger has wandered in, not the way a person checks a house when they believe a demon is haunting them.

He picks up his phone. He puts it down. He stands in the kitchen for a moment. Then he locks the back door, which he has not done once in five weeks, and he goes back to bed.

I am in the living room. The form has a field for ``direct-perception event'' and a field for ``subject response.'' I do not fill them in.
